% !TEX root = ../../main.tex
% C1.3 — Phạm vi của đề tài
% Chương 1: KHÔNG nêu tên công nghệ, mô hình, bộ dữ liệu. RTSP được nêu vì là giao diện đầu vào của camera.
% Phân biệt với 1.4: mục này nói hệ thống LÀM GÌ / KHÔNG LÀM GÌ; mục 1.4 nói những HẠN CHẾ của kết quả đạt được.
\section{Phạm vi của đề tài}
\label{sec:1.3}

Phạm vi của đề tài như sau.

\paragraph{Đối tượng tìm kiếm.} Hệ thống chỉ hỗ trợ tìm kiếm người. Việc tìm kiếm dựa trên đặc điểm ngoại hình tổng thể như trang phục, màu sắc và vật dụng mang theo, không dựa trên khuôn mặt hay các đặc điểm sinh trắc học.

\paragraph{Dữ liệu camera.} Hệ thống nhận video camera theo giao thức truyền phát thời gian thực RTSP (Real Time Streaming Protocol), giao thức phổ biến của camera giám sát. Camera được giả lập bằng cách phát video của một bộ dữ liệu công khai gồm bảy camera tĩnh cùng quan sát một địa điểm; mỗi luồng là một camera được gán vào một khu vực giám sát, và về nguyên tắc thay được bằng camera thật. Các luồng được xử lý lần lượt và kết quả được lưu để tra cứu. Vì máy triển khai chỉ có CPU, không phân tích kịp video độ phân giải cao theo thời gian thực, dữ liệu trình diễn được lập chỉ mục trước từ các tệp video của bộ dữ liệu qua cùng quy trình.

\paragraph{Hình thức mô tả người cần tìm.} Người dùng mô tả người cần tìm bằng một trong ba hình thức:
\begin{itemize}
    \item hình ảnh tham chiếu đã được cắt sẵn, chứa một người;
    \item đoạn mô tả bằng văn bản tiếng Anh;
    \item tập thuộc tính ngoại hình chọn sẵn, gồm giới tính, loại và màu trang phục phía trên, loại và màu trang phục phía dưới, và vật mang theo như ba lô, túi xách.
\end{itemize}
Hệ thống không hỗ trợ truy vấn bằng tiếng Việt và không tự động dịch nội dung mô tả.

\paragraph{Kết quả tìm kiếm.} Mỗi kết quả là một lần xuất hiện của một người trên một camera, xếp theo mức độ phù hợp. Người dùng chọn số kết quả và có thể giới hạn camera, khoảng thời gian. Điểm phù hợp chỉ dùng để xếp hạng trong một lượt tìm; hệ thống không kết luận về danh tính, không lưu lịch sử tìm kiếm, và chỉ ghi vào hồ sơ vụ việc những kết quả người dùng chủ động lưu.

\paragraph{Người dùng.} Ứng dụng phục vụ nhiều nhóm người dùng với quyền hạn khác nhau, được xác định ở Chương~\ref{chapter:phantich}. Việc quản lý danh mục khu vực giám sát không thuộc phạm vi đề tài; danh mục này được khai báo sẵn khi triển khai.

\paragraph{Môi trường triển khai.} Ứng dụng web trên một máy tính cá nhân không có card đồ họa rời, phục vụ trình diễn, không nhằm triển khai thương mại.

\paragraph{Ngoài phạm vi.} Đề tài không bao gồm các nội dung sau:
\begin{itemize}
    \item nhận dạng khuôn mặt hoặc xác định danh tính của người trong hình ảnh;
    \item tìm kiếm phương tiện hoặc các loại đối tượng khác ngoài người;
    \item huấn luyện hoặc tinh chỉnh mô hình trí tuệ nhân tạo;
    \item liên kết các lần xuất hiện của cùng một người giữa nhiều camera khác nhau;
    \item xử lý đồng thời nhiều luồng camera và phát cảnh báo theo thời gian thực.
\end{itemize}
